% ================================ RESUMO =====================================
\begin{center}
\textbf{\sffamily Resumo}
\end{center}
\vspace{-0.85em}
\begin{quote}
\small
O Escore de Risco de Mayo, um modelo de Cox com cinco variáveis de rotina publicado em 1989, ainda é a referência clássica para estimar o prognóstico na cirrose biliar primária, hoje chamada colangite biliar primária. Este estudo mede o valor incremental do aprendizado de máquina sobre uma regressão logística com as mesmas cinco variáveis na predição de óbito em até cinco anos, usando os 312 pacientes rotulados do conjunto público da Mayo Clinic (UCI 878), dos quais 115 morreram no horizonte. O desenho 2 x 2 cruza a família de métodos (regressão e aprendizado de máquina, com \textit{random forest} e XGBoost) e o conjunto de variáveis (5 e 17), com validação cruzada estratificada repetida 5 x 5, ajuste de hiperparâmetros aninhado, validação quase externa nos pacientes não randomizados e um braço de controle de vazamento. A regressão com cinco variáveis obteve AUROC média de 0,906 e inclinação de calibração de 1,04. Nenhum modelo de aprendizado de máquina mostrou ganho: as diferenças de AUROC em relação à referência ficaram entre $-$0,013 e 0,000, todas com intervalos de Bonferroni que incluem zero, e a maior diferença nominal foi a da regressão penalizada com 17 variáveis (+0,009). Usar o tempo de seguimento como entrada, como na competição Kaggle Playground S3E26, reduziu o \textit{log loss} de 0,753 para 0,645 com uma informação indisponível no momento da predição. Neste estudo de caso, o aprendizado de máquina não justificou sua complexidade adicional diante de um modelo de regressão de cinco variáveis.
\end{quote}
\vspace{-0.35em}
\noindent\textbf{Palavras-chave:} colangite biliar primária; escore de Mayo; aprendizado de máquina; valor incremental; calibração; vazamento de dados.
\vspace{0.55em}
